\documentclass[12pt]{article}

\usepackage{latexsym}
\usepackage[T1]{fontenc} 
\usepackage[utf8]{inputenc} 
\usepackage[spanish]{babel}
\usepackage{amsthm, amssymb, amsmath, amsfonts, bbm, mathtools} 
\usepackage{multirow, array}
\usepackage{stmaryrd}
\usepackage{xcolor}
\decimalpoint

\usepackage{wrapfig}
\usepackage{subcaption}

% no tocar nada de aqui
\providecommand{\abs}[1]{\lvert#1\rvert}
\usepackage[hidelinks]{hyperref}
\usepackage{enumerate} 
\usepackage{graphicx}

% entornos personalizados
\newtheorem{example}{Ejemplo}
\newtheorem{notacion}{Notación}
\newtheorem{defi}{Definición}
\newtheorem{obse}{Observación}
\newtheorem{lem}{Lema}

\newcommand{\sol}{\textit{Solución. }}
\renewcommand{\qed}{\hfill $\blacksquare$}
\newcommand{\eej}{\hfill $\square$}

%--------------------- Comandos personalizados ---------------------

% Simbolos de los Numeros
\newcommand{\D}{\mathfrak{D}}
\newcommand{\A}{\mathcal{A}}
\newcommand{\C}{\mathbb{C}}
\newcommand{\F}{\mathbb{F}}
\newcommand{\I}{\mathbb{I}}
\newcommand{\N}{\mathbb{N}}
\newcommand{\Q}{\mathbb{Q}}
\newcommand{\R}{\mathbb{R}}
\newcommand{\Z}{\mathbb{Z}}

% Signos de agrupacion
\newcommand{\pp}[1]{\left( #1 \right)}
\newcommand{\cc}[1]{\left[ #1 \right]}
\newcommand{\set}[1]{\left\{#1\right\}}

\newcommand{\ol}[1]{\overline{#1}}
\newcommand{\ceil}[1]{\left\lceil#1\right\rceil}
\newcommand{\floor}[1]{\left\lfloor#1\right\rfloor}
\newcommand{\ind}[1]{\mathbbm{1}_{\{#1\}}}
\newcommand{\norm}[1]{\left\Vert #1 \right\Vert}

\newcommand{\Text}[1]{\quad\text{#1}\quad}

\newcommand{\bs}{\backslash}
\renewcommand{\c}{\complement}
\newcommand{\beq}{bilipschitz }

% Flechas
\newcommand{\imp}{\Longrightarrow}
\newcommand{\Imp}{\quad\Longrightarrow\quad}
\newcommand{\ssi}{\Longleftrightarrow}
\newcommand{\Ssi}{\quad\Longleftrightarrow\quad}


\newcommand{\e}{\varepsilon}
\renewcommand{\d}{\delta}
\renewcommand{\O}{\Omega}
\renewcommand{\l}{\lambda}
\renewcommand{\a}{\alpha}
\renewcommand{\b}{\beta}
\renewcommand{\abs}[1]{\left\lvert#1\right\rvert}
%\renewcommand{\qedsymbol}{$\blacksquare$}


\newcommand{\Co}{\mathfrak C}
\newcommand{\T}{\mathbb T}
%\newcommand{\uno}{\mathbb 1}
%\newcommand{\ind}[1]{\mathbbm{1}_{\{#1\}}}
\newcommand{\uno}{\mathbbm{1}}


\renewcommand{\o}{\omega}

\newcommand{\x}{\mathrm x}
\newcommand{\y}{\mathrm y}
\newcommand{\Int}{\mathrm{int}}
\newcommand{\Ext}{\mathrm{ext}}
\newcommand{\diam}{\mathrm{diam}}
\newcommand{\cl}{\mathrm{cl}}
\newcommand{\angles}[2]{\left\langle #1, #2\right\rangle}

\textheight=25cm \textwidth=18cm \topmargin=-2cm \oddsidemargin=-1cm

\setlength{\parskip}{0.8em}  % Espaciado párrafos
\setlength{\parindent}{0pt}  % Sangría

\newcommand{\To}{\to}



\begin{document}
    \pagestyle{empty}
    \begin{center}
        {\textsc{Ecuaciones Diferenciales Parciales II} \\
        \medskip 
        \large{Tarea 2}}
    \end{center}
    
    \bigskip
    
    \begin{enumerate}
        \item Diseñar un algoritmo en Python que resuelva ecuaciones
        diferenciales con el método de Descomposición de Adomian en general.

        \item Con el algoritmo creado resolver los siguientes problemas, compara
        con las soluciones analiticas presentadas por Chen en 2004:
        \begin{enumerate}[(a)]
            \item $y'=x^2+y^2$, $y(0)=0$
            \item $u_t=x^2-(u_x)^2/4$, $u(x,0)=0$
            \item $u_t+u^2u_x=0$, $u(x,0)=3x$
            \item \begin{align*}
                \frac{d^2}{dt^2}u+v=0,&\qquad \frac{d^2}{dt^2}v+u=0,\\
                u(0)=0,&\quad u'(0)=1,\\
                v(0)=0,&\quad v'(0)=-1.
            \end{align*}
            \item \begin{align*}
                \frac{\partial^2}{\partial t^2}u=
                \frac{\partial^2}{\partial x^2}u,&\quad 0<x<\pi,~t>0,\\
                u(0,t)=u(\pi, t)=0,&\quad t>0,\\
                u(x,0)=\sin^3(x),&\quad 0<x<\pi,\\
                u_t(x,0) = \sin(2x),&\quad 0<x<\pi.
            \end{align*}
        \end{enumerate}

        \item Resolver los problemas de Black-Scholes:
        \begin{enumerate}[(a)]
            \item \[
                \begin{cases}
                    rC(t, x) = C_t(t, x) + \frac{1}{2}\sigma^2x^2C_{xx}(t, x)
                    + rxC_x(t, x),\qquad& x > 0,~t\in [0, T],\\
                    C(T, x) = \max(x-K, 0),\qquad&\\
                    C(t, x) = x-Ke^{-r(T-t)},\qquad&\text{cuando }x\to\infty,\\
                    C(t, 0) = 0,\qquad&\forall t > 0.
                \end{cases}
            \]

            \item Consideremos el problema:
            \[
                \begin{cases}
                    U_t(\tau, y) = U_{xx}(\tau, y),\quad& y > 0,
                    \tau\in\cc{0, \tfrac{\sigma^2T}{2}},\\
                    U(0, y) = \max\pp{e^{\tfrac{1}{2}(\gamma+1)y}-e^{\tfrac{1}{2}(\gamma-1)y},0},\\
                    U(\tau, L) = e^{\tfrac{1}{2}(\gamma+1)L+\tfrac{1}{4}(\gamma+1)^2\tau}-e^{\tfrac{1}{2}(\gamma-1)L+\tfrac{1}{4}(\gamma-1)^2\tau},\\
                    U(\tau, 0) = 0,
                \end{cases}
            \]
            con
            \[ \tau=\frac{1}{2}\sigma^2(T-t),\qquad y=\log\pp{\frac{x}{K}},
            \qquad \gamma=\frac{2r}{\sigma^2}. \]
            Recordemos que al Usar MDA no se usan las condiciones en la frontera
            como en a), por lo cual se hace una modificación usando
            \begin{align*}
                U(\tau,y) &= u(\tau,y)+w(\tau,y),\\
                w(\tau,y) &= U(0,y)+\pp{U(0,y)-U(\tau,L)}\pp{\frac{y-y_0}{L-y_0}}
            \end{align*}
            el problema se convierte en uno dónde las condiciones de frontera
            son nulas, en este caso, se puede usar el ADM tradicional que usa
            sólo la condición inicial:
            \[
                \begin{cases}
                    u_t(\tau,y)=u_{xx}(\tau, y)-w_t(\tau,y),\quad y>0,
                    ~\tau\in\cc{0,\tfrac{\sigma^2T}{2}},\\
                    u(0,y) = u_0(y)-w(0,y),\\
                    u(\tau,L)=0,\\
                    u(\tau,0)=0,\quad \forall \tau >0,
                \end{cases}
            \]
            con
            \[ u_0(y) = \max\pp{e^{\tfrac{1}{2}(\gamma+1)y}-e^{\tfrac{1}{2}(\gamma-1)y}, 0}. \]
            Con el ADM obtenemos:
            \[ u(t)\approx\psi_k = \sum_{i=0}^{k-1}u_i(t). \]
            La solución es
            \[ C(t, x)=Ke^{-\tfrac{1}{2}(\gamma+1)x-\tfrac{1}{4}(\gamma-1)^{2}\tau}\pp{\psi_k+w(\tau,y)}. \]
            Usa $r=0,05$, $\sigma=0.317$, $K=20$ y $T=0,25$ (3 meses), compara los resultados.

        \end{enumerate}
    \end{enumerate}
\end{document}
